\documentclass[11pt]{article}
\usepackage[T1]{fontenc}
\usepackage{amsmath,amssymb,booktabs}
\usepackage[margin=1in]{geometry}
\usepackage[hidelinks]{hyperref}
\usepackage{listings,tabularx}
\lstset{basicstyle=\ttfamily\small,breaklines=true,columns=fullflexible,
  literate={≤}{{$\le$}}1 {κ}{{$\kappa$}}1 {μ}{{$\mu$}}1 {ℝ}{{$\mathbb R$}}1 {∈}{{$\in$}}1 {∩}{{$\cap$}}1}
\title{Proof Dossier: Strict Center Separation Implies Disjointness}
\author{Project Koios Formal}
\date{}
\begin{document}
\maketitle
\section{Identity and status}
\begin{tabularx}{\textwidth}{@{}l>{\raggedright\arraybackslash}X@{}}
\toprule
Statement ID & \texttt{BG-THM-023} \\
Status & \texttt{proved} \\
Lean declaration & \path{ProjectKoios.BandGauge.Geometry.disjoint_rankTwoFullEllipses_of_add_lt_sqrt_mul_abs_centerSeparation} \\
Owning module & \path{ProjectKoios/BandGauge/Geometry/RankTwoLoss.lean} \\
\bottomrule
\end{tabularx}

\section{Exact statement}
For $\kappa>0$ and nonnegative radii, if
\[
 \tau_s+r<\sqrt\kappa\,|1-\mu|,
\]
then the full spectral and operator ellipses are disjoint.

\section{Source provenance}
This \texttt{formalizes} the strict-separation direction. The source owner is
\path{ksdft2effmass}. Immutable commit:
\begin{center}\texttt{99a46cb59d0c6f27ada7d709f194996ef9619d9f}\end{center}
Source:
\begin{center}\path{docs/publications/papers/Ragasa2026a/v3/manuscript.tex}\end{center}
Locator: Eq. \texttt{exact-intersection}. Later owner worktree changes are excluded.

\section{Proof}
A common point would imply the opposite non-strict inequality by BG-THM-022. The strict
separation premise therefore excludes every common point.

\begin{lstlisting}
theorem disjoint_rankTwoFullEllipses_of_add_lt_sqrt_mul_abs_centerSeparation
    {κ μ spectralBudget operatorRadius : ℝ}
    (hκ : 0 < κ) (hSpectralBudget : 0 ≤ spectralBudget)
    (hOperatorRadius : 0 ≤ operatorRadius)
    (hseparated : spectralBudget + operatorRadius <
      Real.sqrt κ * |1 - μ|) :
    Disjoint (rankTwoFullSpectralEllipse κ spectralBudget)
      (rankTwoFullOperatorEllipse κ μ operatorRadius)
\end{lstlisting}

\section{Boundary and trusted base}
The theorem is classical and not novel. It establishes full-ellipse disjointness, not an
exact positive distance or a rectangle-clipped claim. Lean \texttt{v4.35.0} and pinned
Mathlib report \texttt{[propext, Classical.choice, Quot.sound]}; no admission or
scientific axiom is used.

\section{AI, evidence, and review}
The statement, proof, and dossier were AI-generated or AI-assisted and kernel-checked.
No numerical separation margin or containment premise is validated. No independent
expert review is recorded. On 18 October 2026, repository checks passed. Status:
\texttt{proved}; working draft.
\end{document}
